
\medskip
\noindent\textbf{1.2.3.}\footnote{\label{sga3:VIIB:nde:45}The editors
added the following lemma and introduced the numbering \SGARefLink{sga3:VIIB:1.2.3.A}{1.2.3.A} through
\SGARefLink{sga3:VIIB:1.2.3.F}{1.2.3.F} for later references.}
\phantomsection
\label{sga3:VIIB:1.2.3}
We begin with the following lemma.

\medskip
\noindent\textbf{Lemma 1.2.3.A.}
\phantomsection
\label{sga3:VIIB:1.2.3.A}
{\itshape Let \(k\longrightarrow K\) be a morphism of pseudocompact
rings, let \(B\) be a pseudocompact \(K\)-module, and let \(M\) be a
projective \(k\)-module, without topology. There is a canonical
isomorphism of pseudocompact \(K\)-modules
\begin{equation}
  \theta:
  \operatorname{Hom}_k(M,k)\widehat{\otimes}_k B
  \xrightarrow{\ \sim\ }
  \operatorname{Hom}_k(M,B).
  \tag{2}
  \label{sga3:VIIB:1.2.3:eq2}
\end{equation}}

Here \(M^*=\operatorname{Hom}_k(M,k)\) is endowed with the topology
defined in \SGARefLink{sga3:VIIB:0.2.2}{0.2.2}, which makes it a pseudocompact \(k\)-module. The
topology on \(\operatorname{Hom}_k(M,B)\) is defined analogously: a
basis of neighborhoods of \(0\) is formed by the following
\(K\)-submodules, where \(x\in M\) and \(B'\) is an open
\(K\)-submodule of \(B\):
\[
  \mathscr H(x,B')
  =
  \{f\in\operatorname{Hom}_k(M,B)\mid f(x)\in B'\}.
\]
Finally, \(M^*\widehat{\otimes}_kB\) is the pseudocompact
\(K\)-module obtained from \(M^*\) by base change; cf.\ \SGARefLink{sga3:VIIB:0.5.A}{0.5.A}.

With these conventions, \(\theta\) is plainly an isomorphism when
\(M=k\). Moreover, both sides of \textup{(2)}, regarded as functors in
\(M\), carry direct sums to products; in particular, they commute
with finite direct sums. It follows that \textup{\SGARefLink{sga3:VIIB:1.2.3:eq2}{(2)}} is an
isomorphism first when \(M\) is a free \(k\)-module, and then when
\(M\) is projective.

\medskip
\noindent\textbf{Definition 1.2.3.B.}
\phantomsection
\label{sga3:VIIB:1.2.3.B}
Let \(N\) be a pseudocompact \(k\)-module. We write
\(\mathbf V_k^f(N)\), or
\(\mathbf N^\dagger\), for the
\(\mathbf O_k\)-module which assigns to every
\(C\in\operatorname{Ob}(\operatorname{Alf}/k)\) the \(C\)-module
\[
  (C\widehat{\otimes}_kN)^\dagger,
\]
the dual of \(C\widehat{\otimes}_kN\) in the sense of \SGARefLink{sga3:VIIB:0.2.2}{0.2.2}; that is,
the \(C\)-module
\[
  \operatorname{Hom}_c(C\widehat{\otimes}_kN,C)
  \simeq
  \operatorname{Hom}_c(N,C|_k),
\]
where \(C|_k\) denotes the \(k\)-module obtained from \(C\) by
restriction of scalars. This \(\mathbf O_k\)-module
\(\mathbf V_k^f(N)\) is called the \emph{dual} of \(N\).

If \(k_{\mathfrak m}\) is a local component of \(k\), then for every
object \(C\) of \(\operatorname{Alf}/k\),
\[
  \operatorname{Hom}_c(k_{\mathfrak m},C|_k)
  =
  C_{\mathfrak m}
  =
  C\otimes_k k_{\mathfrak m}
\]
is a projective \(C\)-module. Moreover,
\[
  D\otimes_C C_{\mathfrak m}=D_{\mathfrak m}
\]
for every morphism \(C\longrightarrow D\) in
\(\operatorname{Alf}/k\). Thus
\(\mathbf V_k^f(k_{\mathfrak m})\) is a flat
\(\mathbf O_k\)-module; cf.\ \SGARefLink{sga3:VIIB:1.2.1}{1.2.1}. Since every projective
pseudocompact \(k\)-module is a product of modules
\(k_{\mathfrak m}\) (\SGARefLink{sga3:VIIB:0.2.1}{Proposition~0.2.1}), \SGARefLink{sga3:VIIB:0.2.F}{Corollary~0.2.F} gives:

\medskip
\noindent\textbf{Corollary 1.2.3.C.}
\phantomsection
\label{sga3:VIIB:1.2.3.C}
\textit{The \(\mathbf O_k\)-module \(\mathbf V_k^f(N)\) is flat
whenever \(N\) is a projective pseudocompact \(k\)-module.}

\medskip
\noindent\textbf{Definition 1.2.3.D.}
\phantomsection
\label{sga3:VIIB:1.2.3.D}
Conversely, let \(\mathbf M\) be an admissible
\(\mathbf O_k\)-module (cf.\ \SGARefLink{sga3:VIIB:1.2.1}{1.2.1}).
We call the following pseudocompact \(k\)-module the \emph{dual} of
\(\mathbf M\), and denote it by \(\Gamma^*(\mathbf M)\). As
\(\mathfrak l\) ranges over the open ideals of \(k\), endow each
\(k/\mathfrak l\)-module
\[
  \mathbf M(k/\mathfrak l)^*
  =
  \operatorname{Hom}_{k/\mathfrak l}
    \bigl(\mathbf M(k/\mathfrak l),k/\mathfrak l\bigr)
\]
with the topology described in \SGARefLink{sga3:VIIB:0.2.2}{0.2.2}; it is then a pseudocompact
\(k\)-module. Since
\[
  \mathbf M(k/\mathfrak l)
  \simeq
  \mathbf M(k/\mathfrak l')
    \otimes_{k/\mathfrak l'}(k/\mathfrak l)
  \qquad
  (\mathfrak l\supset\mathfrak l'),
\]
there are transition morphisms
\[
  \begin{split}
  \operatorname{Hom}_{k/\mathfrak l'}
    \bigl(\mathbf M(k/\mathfrak l'),k/\mathfrak l'\bigr)
  &\longrightarrow
  \operatorname{Hom}_{k/\mathfrak l'}
    \bigl(\mathbf M(k/\mathfrak l'),k/\mathfrak l\bigr)\\
  &=
  \operatorname{Hom}_{k/\mathfrak l}
    \bigl(\mathbf M(k/\mathfrak l),k/\mathfrak l\bigr).
  \end{split}
\]
By definition, \(\Gamma^*(\mathbf M)\) is the projective limit of this
filtered projective system of pseudocompact \(k\)-modules.

From now on, suppose in addition that \(\mathbf M\) is a flat
\(\mathbf O_k\)-module (cf.\ \SGARefLink{sga3:VIIB:1.2.1}{1.2.1}). Each
\(\mathbf M(k/\mathfrak l')\) is then a projective
\(k/\mathfrak l'\)-module, so the transition morphisms above are
surjective. Consequently the projections
\[
  \Gamma^*(\mathbf M)
  \longrightarrow
  \mathbf M(k/\mathfrak l)^*
\]
are also surjective, since filtered projective limits are exact in
\(\operatorname{PC}(k)\).

Let \(\tau:k\longrightarrow K\) be a morphism of pseudocompact rings.
We write \(\mathbf M\otimes_kK\), or simply \(\mathbf M_K\), for the
\(\mathbf O_K\)-functor defined as follows. If \(C\) is a
finite-length \(K\)-algebra, then the kernel of
\(k\longrightarrow C\) is an open ideal \(\mathfrak l\), and we set
\[
  \mathbf M_K(C)
  =
  \mathbf M(k/\mathfrak l)\otimes_k C.
\]
The same value is obtained from
\(\mathbf M(k/\mathfrak l')\otimes_kC\) for every open ideal
\(\mathfrak l'\subset\mathfrak l\). We then define
\(\Gamma_K^*(\mathbf M_K)\) as the projective limit, with \(I\)
ranging over the open ideals of \(K\), of the pseudocompact
\(K\)-modules
\[
  \operatorname{Hom}_{K/I}\bigl(\mathbf M_K(K/I),K/I\bigr)
  =
  \operatorname{Hom}_{k/\mathfrak l}
    \bigl(\mathbf M(k/\mathfrak l),K/I\bigr),
\]
where on the right \(\mathfrak l\) is any open ideal of \(k\) such
that \(\tau(\mathfrak l)\subset I\).

By \SGARefLink{sga3:VIIB:1.2.3.A}{Lemma~1.2.3.A}, the right-hand term identifies with
\[
  \mathbf M(k/\mathfrak l)^*
  \widehat{\otimes}_k(K/I).
\]
Since the projections
\(\Gamma^*(\mathbf M)\longrightarrow\mathbf M(k/\mathfrak l)^*\)
are surjective, the projective limit of these latter modules is
precisely the pseudocompact \(K\)-module
\[
  \Gamma^*(\mathbf M)\widehat{\otimes}_kK
\]
(cf.\ \SGARefLink{sga3:VIIB:0.5.A}{0.5.A}). We have therefore proved that, for every flat
\(\mathbf O_k\)-module \(\mathbf M\), formation of
\(\Gamma^*(\mathbf M)\) commutes with extension of the base; that is,
\begin{equation}
  \Gamma_K^*(\mathbf M\otimes_kK)
  \simeq
  \Gamma^*(\mathbf M)\widehat{\otimes}_kK.
  \tag{\(\star\)}
  \label{sga3:VIIB:1.2.3:star}
\end{equation}

\medskip
\noindent\textbf{Proposition 1.2.3.E.}
\phantomsection
\label{sga3:VIIB:1.2.3.E}
{\itshape
\begin{enumerate}
  \item[\textup{(i)}]
  The functors
  \[
    N\longmapsto\mathbf V_k^f(N),
    \qquad
    \mathbf M\longmapsto\Gamma^*(\mathbf M)
  \]
  define an anti-equivalence between the category of projective
  pseudocompact \(k\)-modules and the category of flat
  \(\mathbf O_k\)-modules.\footnote{\label{sga3:VIIB:nde:48}The
  editors added assertion \textup{\SGARefLink{sga3:vii:paragraph:5-1:equation:ii}{(ii)}} below and the proof that
  follows. The original said: ``If \(\mathbf M\) is a flat
  \(\mathbf O_k\)-module, it is clear that \(\mathbf M\) is the dual
  of \(\mathbf M^*\); hence the functors \(N\mapsto N^*\) and
  \(\mathbf M\mapsto\mathbf M^*\) define an
  anti-equivalence \ldots''}

  \item[\textup{(ii)}]
  Moreover, if \(k\longrightarrow K\) is a morphism of pseudocompact
  rings, the preceding anti-equivalence ``commutes with base change''
  in the following sense: if
  \(N\simeq\Gamma^*(\mathbf M)\), then
  \[
    N\widehat{\otimes}_kK
    \simeq
    \Gamma_K^*(\mathbf M\otimes_kK).
  \]
\end{enumerate}}

\noindent\textit{Proof.}
First, there is a natural transformation
\[
  \phi_N:
  N\longrightarrow\Gamma^*\bigl(\mathbf V_k^f(N)\bigr).
\]
The functor \(\Gamma^*\circ\mathbf V_k^f\) commutes with products by
\SGARefLink{sga3:VIIB:0.3.5}{0.3.5} and \SGARefLink{sga3:VIIB:0.2.F}{0.2.F}, so it is enough to verify that \(\phi_N\) is an
isomorphism when \(N\) is a local component \(k_{\mathfrak m}\) of
\(k\). In this case, for every open ideal \(\mathfrak l\) of \(k\)
contained in \(\mathfrak m\), the natural morphism
\[
  (k/\mathfrak l)_{\mathfrak m}
  \longrightarrow
  \operatorname{Hom}_{k/\mathfrak l}
  \left(
    \operatorname{Hom}_c
      (k_{\mathfrak m}\widehat{\otimes}_k k/\mathfrak l,
       k/\mathfrak l),
    k/\mathfrak l
  \right)
\]
is an isomorphism, which proves the assertion.

Conversely, let \(\mathbf M\) be a flat \(\mathbf O_k\)-module. We
show that
\[
  \Gamma^*(\mathbf M)
  =
  \varprojlim_{\mathfrak l}\mathbf M(k/\mathfrak l)^*
\]
is a projective object of \(\operatorname{PC}(k)\). Let
\(N\longrightarrow N'\) be a surjective morphism between objects of
\(\operatorname{LF}(k)\). By \SGARefLink{sga3:VIIB:0.2.F}{0.2.F}\,\textup{\SGARefLink{sga3:vii:paragraph:5-1:equation:i}{(i)}} and
\textup{\SGARefLink{sga3:vii:paragraph:5-1:equation:ii}{(ii)}}, it is enough to prove that the natural map
\[
  \varinjlim_{\mathfrak l}
    \operatorname{Hom}_c(\mathbf M(k/\mathfrak l)^*,N)
  \longrightarrow
  \varinjlim_{\mathfrak l}
    \operatorname{Hom}_c(\mathbf M(k/\mathfrak l)^*,N')
\]
is surjective. This is clear: \(N\) and \(N'\) are
\(k/\mathfrak l_0\)-modules for some open ideal
\(\mathfrak l_0\), and if
\(\mathfrak l\subset\mathfrak l_0\), every morphism
\(\mathbf M(k/\mathfrak l)^*\longrightarrow N'\) lifts to a morphism
\(\mathbf M(k/\mathfrak l)^*\longrightarrow N\), since
\(\mathbf M(k/\mathfrak l)^*\) is a projective object of
\(\operatorname{PC}(k/\mathfrak l)\).

There is a morphism
\[
  \psi:
  \mathbf M\longrightarrow
  \mathbf V_k^f\bigl(\Gamma^*(\mathbf M)\bigr)
\]
of functors from \(\operatorname{Alf}/k\) to
\((\mathbf{Sets})\). Let \(B\) be an object of
\(\operatorname{Alf}/k\). We claim that
\begin{equation}
  \begin{split}
  \psi(B):\mathbf M(B)
  &\longrightarrow
  \mathbf V_k^f\bigl(\Gamma^*(\mathbf M)\bigr)(B)\\
  &=
  \operatorname{Hom}_c
  \left(
    \varprojlim_{\mathfrak l}
      \bigl(\mathbf M(k/\mathfrak l)^*
        \widehat{\otimes}_kB\bigr),
    B
  \right)
  \end{split}
  \tag{1}
  \label{sga3:VIIB:1.2.3:eq1}
\end{equation}
is a bijection. In the equality above we used the fact that
\(\widehat{\otimes}\) commutes with filtered projective limits. Let
\(\mathfrak l_0\) be an open ideal of \(k\) contained in the kernel
of \(k\longrightarrow B\). By \SGARefLink{sga3:VIIB:1.2.3.A}{Lemma~1.2.3.A}, for every
\(\mathfrak l\subset\mathfrak l_0\), there is a canonical
isomorphism of pseudocompact \(B\)-modules
\[
  \mathbf M(k/\mathfrak l)^*\widehat{\otimes}_kB
  \xrightarrow{\ \sim\ }
  \operatorname{Hom}_{k/\mathfrak l}
    \bigl(\mathbf M(k/\mathfrak l),B\bigr).
\]
Since
\[
  \mathbf M(B)
  =
  \mathbf M(k/\mathfrak l)
    \otimes_{k/\mathfrak l}B,
\]
the right-hand side is
\(\operatorname{Hom}_B(\mathbf M(B),B)\). Thus the projective system
in \textup{\SGARefLink{sga3:VIIB:1.2.3:eq1}{(1)}} is constant for
\(\mathfrak l\subset\mathfrak l_0\), and \textup{(1)} reduces to the
canonical morphism
\[
  \mathbf M(B)
  \longrightarrow
  \operatorname{Hom}_c
  \bigl(\operatorname{Hom}_B(\mathbf M(B),B),B\bigr).
\]
This is an isomorphism by \SGARefLink{sga3:VIIB:0.2.2}{0.2.2}, since \(B\) is Artinian and
\(\mathbf M(B)\) is a projective \(B\)-module. This proves assertion
\textup{\SGARefLink{sga3:vii:paragraph:5-1:equation:i}{(i)}} of the proposition; assertion \textup{\SGARefLink{sga3:vii:paragraph:5-1:equation:ii}{(ii)}} follows
from the isomorphism \textup{(\(\star\))} established above.
\hfill\(\square\)

\medskip
\noindent\textbf{Remark 1.2.3.F.}
\phantomsection
\label{sga3:VIIB:1.2.3.F}
Let us return to the statement of the proposition and suppose, in
addition, that \(N\) is a topologically free pseudocompact
\(k\)-module. In this case one can choose, ``coherently,'' bases for
the \(C\)-modules \(\mathbf V_k^f(N)(C)\).

Indeed, let \((n_i)\) be a pseudobasis of \(N\) (\SGARefLink{sga3:VIIB:0.2.1}{0.2.1}), and let
\(n_i^C\) be the canonical image of \(n_i\) in
\(C\widehat{\otimes}_kN\), for
\(C\in\operatorname{Alf}/k\). Define
\(\delta_i^C\in(C\widehat{\otimes}_kN)^\dagger\) by
\[
  \delta_i^C(n_i^C)=1,
  \qquad
  \delta_i^C(n_j^C)=0\quad(i\ne j).
\]
Then \((\delta_i^C)\) is a basis of
\(\mathbf V_k^f(N)(C)\). Moreover, for every morphism
\(\varphi:C\longrightarrow D\) in \(\operatorname{Alf}/k\), the map
\(\mathbf V_k^f(N)(\varphi)\) sends \(\delta_i^C\) to
\(\delta_i^D\).

